\documentclass[11pt,a4paper]{ctexart}
\usepackage[teacher]{ipara}

\pagestyle{fancy}
\fancyhf{}
\lhead{\small 专题讲义：函数的对称性与周期性（教师备课版）}
\rhead{\small 高中数学一对一}
\cfoot{\small 第 \thepage\ 页\quad 共 \zpageref{LastPage} 页}
\renewcommand{\headrulewidth}{0.4pt}
\renewcommand{\footrulewidth}{0pt}

\newcommand{\lessoninfo}[2]{%
  \begin{center}
    {\LARGE\bfseries #1 \quad (教师版)}\par
    \vspace{0.8em}
    {\small 学生：#2 \quad 时长：120 分钟 \quad 讲次：专题突破}
  \end{center}
  \vspace{0.6em}
}

\begin{document}

\lessoninfo{函数的对称性与周期性}{\blank{2.4cm}}

\section{备课指导与学情预判}

\begin{tabularx}{\textwidth}{p{2.8cm}p{3.2cm}X}
\toprule
时间 & 环节 & 教师动作与板书要点\\
\midrule
00--15 分钟 & 几何图示引入 & 画出不规则波浪线，通过 3 组 TikZ 图演示翻折与旋转过程。主打“上图下式”板书，让学生看见平移。\\ \hline
15--45 分钟 & 三大模型代数推导 & 示范箭号流程 $x \to 2a-x \to x+2(b-a)$。重点说明两中心**等高**前提，以及一轴一中心**反相平移（两步恢复）**。\\ \hline
45--65 分钟 & 条件反射与口诀 & 归纳 4 个核心结构与总口诀，提问抽查学生对符号 $f(a+x)=f(b-x)$ 与 $f(a+x)=f(x+b)$ 的辨析。\\ \hline
65--90 分钟 & 降维五大应用与秒杀 & 讲解大数化小、零点复制、积分抵消及联手秒杀性质（严格单调、极限、有界性）。\\ \hline
90--120 分钟 & 四步法演练与复盘 & 指导学生写出规范的“固定四步法”步骤，纠正“凭空换元”与“盲言最小正周期”问题。\\
\bottomrule
\end{tabularx}

\vspace{0.8em}

\begin{teachbox}{核心教学指导思想}
学生卡住的地方往往不是公式，而是**没有看见：图像翻折两次之后，为什么会变成平移。**\par
教学切忌一上来就写 $f(2a-x)=f(x), f(2b-x)=f(x)$ 强行换元，否则基础弱的学生会觉得是在凭空捏造。最理想的讲法是：\textbf{先让学生在图上看见平移，再用代数把这个平移证明出来。}
\end{teachbox}

\newpage

\section{核心考向 1：几何直观与“上图下式”板书结构}

\teachblock{上图下式板书规范}{
教师在黑板上画图时，必须保持“上面画图形变换、中间写横坐标轨迹、下面写函数关系”的三层对应关系。
}

\vspace{0.5em}

\teacherproblem{一}{两条对称轴复合：轴对称 + 轴对称 = 平移}

\teachvalue{直观展示为何两平行对称轴翻折两次后，曲线的形状方向完全恢复正立，仅产生 $2(b-a)$ 的水平平移。}

\begin{center}
\begin{tikzpicture}[scale=0.85, >=stealth]
  \draw[->, gray!60, thin] (-0.5,0) -- (7.5,0) node[right, black] {$x$};
  \draw[->, gray!60, thin] (0,-0.8) -- (0,2.5) node[above, black] {$y$};
  \draw[dashed, thick] (2,-0.6) -- (2,2.3) node[above] {$x=a$};
  \draw[dashed, thick] (4,-0.6) -- (4,2.3) node[above] {$x=b$};
  
  % C0: x in [0,2]
  \draw[very thick] plot[smooth, tension=0.7] coordinates {(0, 0.5) (1, 1.9) (2, 0.8)};
  \node[above] at (1, 1.9) {$C_0$};
  \fill (1, 1.9) circle (1.5pt) node[above left] {$P(x,y)$};
  
  % C1: x in [2,4] (关于 x=a=2 轴对称)
  \draw[thick, densely dotted] plot[smooth, tension=0.7] coordinates {(2, 0.8) (3, 1.9) (4, 0.5)};
  \node[above] at (3, 1.9) {$C_1$};
  \fill (3, 1.9) circle (1.5pt) node[above right] {$P_1(2a-x, y)$};
  \draw[->, bend left=25, darkgray] (1.0, 2.1) to node[above, font=\small] {$x=a$ 翻折} (3.0, 2.1);
  
  % C2: x in [4,6] (关于 x=b=4 轴对称)
  \draw[very thick] plot[smooth, tension=0.7] coordinates {(4, 0.5) (5, 1.9) (6, 0.8)};
  \node[above] at (5, 1.9) {$C_2$};
  \fill (5, 1.9) circle (1.5pt) node[above right] {$P_2(x+2(b-a), y)$};
  \draw[->, bend left=25, darkgray] (3.0, 2.1) to node[above, font=\small] {$x=b$ 翻折} (5.0, 2.1);
  
  \draw[->, double, thick] (1.0, -0.4) -- (5.0, -0.4) node[midway, below, font=\small\bfseries] {无缝连贯：整体向右平移 $2(b-a)$};
\end{tikzpicture}
\end{center}

\begin{paracol}{2}
\teachblock{标准代数证明}{
设 $P(x,y)$ 为 $f(x)$ 图像上任意一点。\par
1. 关于 $x=a$ 轴对称得点 $P_1(2a-x, y)$，故 $f(2a-x) = f(x)$；\par
2. 再关于 $x=b$ 轴对称得点 $P_2\bigl(2b-(2a-x), y\bigr) = P_2\bigl(x+2(b-a), y\bigr)$；\par
3. 连续代入得：
\[
\begin{aligned}
f\bigl(x+2(b-a)\bigr) &= f\bigl(2b-(2a-x)\bigr) \\
&= f(2a-x) \\
&= f(x).
\end{aligned}
\]
故 $2|b-a|$ 是 $f(x)$ 的一个周期。
}

\switchcolumn

\teachblock{学生问题与诊断}{
【学生问题】\par
学生常问：“为什么第一次是 $2a-x$，第二次变成 $2b-(2a-x)$？”\par
【课堂处理】\par
用手在图上指明：第二次翻折时，被翻折的点的横坐标已经是 $X = 2a-x$。根据轴对称公式 $X \to 2b-X$，直接把 $X=2a-x$ 整体代入，即可得到 $2b-(2a-x) = x+2(b-a)$。\par
强调：**只算横坐标，不要急着看函数值**。
}

\teachblock{追问设计}{
追问：如果 $b < a$，平移方向是向左还是向右？周期公式 $2|b-a|$ 是否依然成立？（引导学生注意到绝对值的几何意义是“轴间距的 2 倍”）。
}
\end{paracol}

\newpage

\teacherproblem{二}{两个对称中心复合：中心对称 + 中心对称 = 平移}

\teachvalue{帮助学生理清为什么两个对称中心**必须等高**才能产生周期；若不等高，则产生“斜向重复”的准周期。}

\begin{center}
\begin{tikzpicture}[scale=0.85, >=stealth]
  \draw[->, gray!60, thin] (-0.5,0) -- (7.5,0) node[right, black] {$x$};
  \draw[->, gray!60, thin] (0,-0.8) -- (0,2.5) node[above, black] {$y$};
  \draw[gray!40, dashed] (-0.5, 0.8) -- (7.5, 0.8) node[right, black, font=\small] {$y=m$};
  
  \fill (2.0, 0.8) circle (2pt) node[below right] {$A(a,m)$};
  \fill (4.0, 0.8) circle (2pt) node[below right] {$B(b,m)$};
  
  % C0: x in [0,2]
  \draw[very thick] plot[smooth, tension=0.7] coordinates {(0, 0.8) (1, 1.9) (2, 0.8)};
  \node[above] at (1, 1.9) {$C_0$};
  \fill (1, 1.9) circle (1.5pt) node[above left] {$P(x,y)$};
  
  % C1: x in [2,4] (绕 A(2,0.8) 旋转 180 度)
  \draw[thick, densely dotted] plot[smooth, tension=0.7] coordinates {(2, 0.8) (3, -0.3) (4, 0.8)};
  \node[below] at (3, -0.3) {$C_1$};
  \fill (3, -0.3) circle (1.5pt) node[below right] {$P_1(2a-x, 2m-y)$};
  \draw[dashed, gray] (1, 1.9) -- (3, -0.3);
  
  % C2: x in [4,6] (绕 B(4,0.8) 旋转 180 度)
  \draw[very thick] plot[smooth, tension=0.7] coordinates {(4, 0.8) (5, 1.9) (6, 0.8)};
  \node[above] at (5, 1.9) {$C_2$};
  \fill (5, 1.9) circle (1.5pt) node[above right] {$P_2(x+2(b-a), y)$};
  \draw[dashed, gray] (3, -0.3) -- (5, 1.9);
  
  \draw[->, double, thick] (1.0, 2.2) -- (5.0, 2.2) node[midway, above, font=\small\bfseries] {无缝连贯：中心对称两次 $=$ 向右平移 $2(b-a)$};
\end{tikzpicture}
\end{center}

\begin{paracol}{2}
\teachblock{标准推导与不等高拓展}{
设两点为 $A(a,m)$ 与 $B(b,n)$：\par
1. 关于 $A$ 对称：$f(2a-x) = 2m - f(x)$；\par
2. 关于 $B$ 对称：$f(2b-x) = 2n - f(x)$；\par
3. 连续复合代入：
\[
\begin{aligned}
f\bigl(x+2(b-a)\bigr) &= f\bigl(2b-(2a-x)\bigr) \\
&= 2n - f(2a-x) \\
&= 2n - \bigl(2m - f(x)\bigr) \\
&= f(x) + 2(n-m).
\end{aligned}
\]
【分支讨论】：
\begin{itemize}[leftmargin=1.2em, itemsep=0.2em]
  \item \textbf{若 $m=n$}：$f\bigl(x+2(b-a)\bigr) = f(x) \implies 2|b-a|$ 是周期；
  \item \textbf{若 $m \neq n$}：$f(x+T) = f(x) + C \implies$ 每平移一段，高度上升或下降 $2(n-m)$（准周期/阶梯上升）。
\end{itemize}
}

\switchcolumn

\teachblock{学生错因与反思}{
【学生问题】\par
学生常忽略 $m=n$ 这个前提，以为任意两个对称中心都能产生周期。\par
【现场教学点拨】\par
在黑板上画一个不等高的例子：点 $A(1,1)$ 和点 $B(3,2)$。旋转一次后图像翻转颠倒，再旋转一次虽然横坐标恢复平移，但 $y$ 坐标变成了 $y + 2(2-1) = y+2$。图像像台阶一样逐段向上爬升。\par
明确告诉学生：**只有纵坐标等高的两个中心，才能消去 $y$ 方向的变化，形成纯粹的水平平移！**
}
\end{paracol}

\newpage

\teacherproblem{三}{轴对称 + 中心对称：反相平移（两步恢复，$T=4|b-a|$）}

\teachvalue{解释为什么一轴一中心的周期是 $4|b-a|$ 而不是 $2|b-a|$，突出“两步恢复”的原理。}

\begin{center}
\begin{tikzpicture}[scale=0.8, >=stealth]
  \draw[->, gray!60, thin] (-0.5,0) -- (10.5,0) node[right, black] {$x$};
  \draw[->, gray!60, thin] (0,-1.5) -- (0,2.5) node[above, black] {$y$};
  \draw[dashed, thick] (2,-1.2) -- (2,2.3) node[above] {$x=a$};
  \draw[gray!40, dashed] (-0.5, 0.8) -- (10.5, 0.8) node[right, black, font=\small] {$y=m$};
  \fill (4.0, 0.8) circle (2pt) node[below right] {$B(b,m)$};
  
  % C0: x in [0,2]
  \draw[very thick] plot[smooth, tension=0.7] coordinates {(0, 0.8) (1, 2.0) (2, 0.8)};
  \node[above] at (1, 2.0) {$C_0$ (正位)};
  
  % C1: x in [2,4]
  \draw[thick, densely dotted] plot[smooth, tension=0.7] coordinates {(2, 0.8) (3, 2.0) (4, 0.8)};
  \node[above] at (3, 2.0) {$C_1$};
  
  % C2: x in [4,6]
  \draw[thick, densely dotted] plot[smooth, tension=0.7] coordinates {(4, 0.8) (5, -0.4) (6, 0.8)};
  \node[below] at (5, -0.4) {$C_2$ (颠倒: $f(x+T)=2m-f(x)$)};
  
  % C3: x in [6,8]
  \draw[thick, densely dotted] plot[smooth, tension=0.7] coordinates {(6, 0.8) (7, -0.4) (8, 0.8)};
  \node[below] at (7, -0.4) {$C_3$};
  
  % C4: x in [8,10]
  \draw[very thick] plot[smooth, tension=0.7] coordinates {(8, 0.8) (9, 2.0) (10, 0.8)};
  \node[above] at (9, 2.0) {$C_4$ (恢复正立)};
  
  \draw[->, double, thick] (1.0, 2.3) -- (9.0, 2.3) node[midway, above, font=\small\bfseries] {无缝连续曲线：两步恢复周期为 $4|b-a|$};
\end{tikzpicture}
\end{center}

\begin{paracol}{2}
\teachblock{代数证明}{
已知 $f(2a-x) = f(x)$ 且 $f(2b-x) = 2m - f(x)$：\par
一轮复合：
\[
\begin{aligned}
f\bigl(x+2(b-a)\bigr) &= f\bigl(2b-(2a-x)\bigr) \\
&= 2m - f(2a-x) \\
&= 2m - f(x). \quad (\text{反相平移})
\end{aligned}
\]
二轮复合（再平移一次 $2(b-a)$）：
\[
\begin{aligned}
f\bigl(x+4(b-a)\bigr) &= 2m - f\bigl(x+2(b-a)\bigr) \\
&= 2m - (2m - f(x)) \\
&= f(x).
\end{aligned}
\]
故 $4|b-a|$ 是 $f(x)$ 的一个周期。
}

\switchcolumn

\teachblock{学生诊断与突破}{
【核心错因】\par
学生算到 $f(x+2(b-a)) = 2m-f(x)$ 时，常常激动地直接下结论说周期是 $2|b-a|$。\par
【教师点拨口诀】\par
“**平移加颠倒，还算不上周期；必须再平移一次颠倒回来，才叫完形归位！**”\par
对比：
\begin{itemize}[leftmargin=1em, itemsep=0.2em]
  \item $f(x+T) = f(x) \implies$ 周期为 $T$；
  \item $f(x+T) = -f(x) \implies$ 周期为 $2T$；
  \item $f(x+T) = 2m - f(x) \implies$ 周期为 $2T$。
\end{itemize}
因此这里 $T_0 = 2|b-a|$，总周期为 $2T_0 = 4|b-a|$。
}
\end{paracol}

\newpage

\section{核心考向 2：条件反射与周期性“降维”五大应用}

\teacherproblem{四}{对称与周期条件反射表达与记忆口诀}

\teachblock{四种表达的总结与辨析}{
\begin{enumerate}[leftmargin=1.5em, itemsep=0.3em]
  \item \textbf{轴对称}：$f(a+x) = f(a-x) \iff f(2a-x) = f(x)$（对称轴 $x=a$）。注意与 $f(a+x)=f(x-a)$ 区分，$x$ 符号相反是对称，符号相同是周期。
  \item \textbf{中心对称}：$f(a+x) + f(a-x) = 2b \iff f(2a-x) = 2b - f(x)$（对称中心 $(a,b)$）。
  \item \textbf{直接周期}：$f(x+T) = f(x)$（周期 $T$）。
  \item \textbf{反周期关系}：$f(x+T) = -f(x)$ 或 $f(x+T) = 2m - f(x)$（周期 $2|T|$）。
\end{enumerate}
}

\begin{teachbox}{总记忆口诀}
\centering\bfseries\large
一轴一翻，两轴一移；两心一移，轴心反相；反相两次，必出周期。
\end{teachbox}

\vspace{0.8em}

\teacherproblem{五}{周期性的“降维打击”与三大联手秒杀定理}

\begin{paracol}{2}
\teachblock{周期性的五大解题功能}{
1. \textbf{无限压缩到有限}：研究整个 $\mathbb{R}$ 转化为研究一个周期 $[0,T)$；\par
2. \textbf{大数化小}：利用 $f(x+nT) = f(x)$ 对求值项取模；\par
3. \textbf{零点成组复制}：一个零点 $x_0$ 派生出无穷解 $x_0 + nT$；\par
4. \textbf{最值单调重复}：只需确定主周期内的增减性与极值；\par
5. \textbf{定积分搬移}：$\int_{a}^{a+T} f(x) dx = \int_{0}^{T} f(x) dx$。
}

\switchcolumn

\teachblock{三大联手秒杀定理}{
【定理 1：周期 + 严格单调】\par
若 $f(x)$ 定义在 $\mathbb{R}$ 上且具有周期性，则 $f(x)$ **绝不可能在 $\mathbb{R}$ 上严格单调**（除非定义域非 $\mathbb{R}$）。\par
【定理 2：周期 + 无穷远极限】\par
若 $\lim_{x\to+\infty} f(x) = A$ 且 $f(x)$ 为周期函数，则 **$f(x) \equiv A$ 必为常数函数**。\par
【定理 3：连续周期函数的有界性】\par
若 $f(x)$ 在 $\mathbb{R}$ 上连续且有周期，则 $f(x)$ **在 $\mathbb{R}$ 上必定有界**。
}
\end{paracol}

\newpage

\section{核心考向 3：典例演练与“固定四步法”示范}

\teacherproblem{例题 1}{轴对称与中心对称复合综合求值}

{\small
已知定义在 $\mathbb{R}$ 上的函数 $f(x)$ 满足 $f(1+x) = f(1-x)$，且 $f(3+x) = -f(3-x)$。
(1) 求证：$f(x)$ 是周期函数，并求出其一个周期；
(2) 若 $f(1) = 2$，求 $f(2025) + f(2026)$ 的值。
}

\teachvalue{训练学生规范使用“固定四步法”，辨析轴对称与中心对称的代数特征，以及大数化小的取模技巧。}

\begin{paracol}{2}
\teachblock{标准解答与板书}{
【第一步：翻译几何条件】\par
由 $f(1+x) = f(1-x)$ 得 $f(x)$ 关于直线 $x=1$ 对称，即：
\[
f(2-x) = f(x). \quad \text{(式 1)}
\]
由 $f(3+x) = -f(3-x)$ 得 $f(x)$ 关于点 $(3,0)$ 中心对称，即：
\[
f(6-x) = -f(x). \quad \text{(式 2)}
\]

【第二步与第三步：轨迹计算与连续代换】\par
计算横坐标两次变换：$x \to 2-x \to 6-(2-x) = x+4$。
\[
\begin{aligned}
f(x+4) &= f\bigl(6-(2-x)\bigr) \\
&= -f(2-x) \quad (\text{代入式 2}) \\
&= -f(x). \quad (\text{代入式 1})
\end{aligned}
\]
再平移一次：$f(x+8) = -f(x+4) = -(-f(x)) = f(x)$。

【第四步：下结论】\par
故 $8$ 是 $f(x)$ 的一个周期。

【第 (2) 问求值】\par
由式 1 得 $f(1) = f(2-1) = f(1) = 2$；\par
由式 2 得 $f(5) = -f(1) = -2$；由 $f(x+4)=-f(x)$ 得 $f(3) = -f(-1) = 0$。\par
对于大数取模：\par
$2025 = 253 \times 8 + 1 \implies f(2025) = f(1) = 2$；\par
$2026 = 253 \times 8 + 2 \implies f(2026) = f(2) = f(2-0) = f(0)$。\par
而由 $f(x+4)=-f(x)$ 得 $f(4)=-f(0)$，又 $f(6-2)=-f(2) \implies f(4)=-f(2) \implies f(2)=0$。\par
故 $f(2025) + f(2026) = 2 + 0 = 2$。
}

\switchcolumn

\teachblock{学生问题与诊断}{
【常见书写漏洞】\par
1. 第一步将 $f(3+x)=-f(3-x)$ 写错为关于 $x=3$ 轴对称；\par
2. 误以为周期是 $4$，忽略了前面的负号 $-f(x)$ 导致周期少算一半；\par
3. 第 (2) 问大数化小对 $f(2)$ 的求值卡壳。

\teachblock{课堂处理与追问}{
追问：如果小明直接记口诀“一轴一中心，周期 $4|b-a|$”，这里 $a=1, b=3$，周期 $4 \times |3-1| = 8$，能否直接写出周期 8？\par
教师指导：高考解答题中，**口诀可以作为快速校验答案的工具，但书写步骤必须按四步法写出代换过程**，否则会扣得分点！
}

\remedy{
若学生大数取模失误，课后安排 2 道周期取模求值变式重练。
}
}
\end{paracol}

\newpage

\teacherproblem{例题 2}{反周期关系与零点成组复制}

{\small
设 $f(x)$ 是定义在 $\mathbb{R}$ 上的连续函数，满足 $f(x+2) = -f(x)$。
(1) 求 $f(x)$ 的一个周期；
(2) 若 $f(x)$ 在 $[0,2]$ 上的零点个数为 $3$，求 $f(x)$ 在 $[0,10]$ 上的零点个数。
}

\teachvalue{训练学生掌握“反周期关系”的周期导出，以及利用周期性进行零点复制时对端点重复问题的精确把控。}

\begin{paracol}{2}
\teachblock{标准解答}{
(1) 由 $f(x+2) = -f(x)$，得：
\[
f(x+4) = -f(x+2) = -(-f(x)) = f(x).
\]
故 $T=4$ 是 $f(x)$ 的一个周期。

(2) 由 $f(x+2) = -f(x)$，令 $x=0$ 得 $f(2) = -f(0)$。\par
又 $f(x)$ 为连续函数且在 $[0,2]$ 上有 3 个零点：\par
假设零点为 $x_1, x_2, x_3 \in [0,2]$。\par
因为 $f(x+2) = -f(x)$，若 $x_0$ 是 $f(x)$ 的零点，则 $f(x_0+2) = -f(x_0) = 0$，即 $x_0+2$ 也是 $f(x)$ 的零点。\par
因此在区间 $[2,4]$ 上，零点恰好为 $x_1+2, x_2+2, x_3+2$，也有 3 个零点。\par
注意端点重合情况：\par
若 $0$ 是零点，则 $2$ 必为零点， $4$ 也必为零点。区间 $[0,2]$ 与 $[2,4]$ 的交点 $x=2$ 被重复计数。\par
在半闭半开主周期区间 $[0,4)$ 内，零点个数为 $3 + 3 - 1 = 5$ 个。\par
区间 $[0,10]$ 的长度为 $10$，包含：\par
2 个完整主周期 $[0,8)$（共 $2 \times 5 = 10$ 个零点），以及剩余区间 $[8,10]$（与 $[0,2]$ 同构，有 3 个零点）。\par
但 $x=8$ 在前面 $[0,8)$ 中已计算，故 $[8,10]$ 贡献额外 3 个零点。\par
总零点个数为 $10 + 3 = 13$ 个。
}

\switchcolumn

\teachblock{学生诊断与警示}{
【易错点：端点重复计算】\par
学生极易简单粗暴写出 $3 \times 5 = 15$ 个，完全忽视了区间拼接处端点零点的重复计数！\par
【教师画图指导】\par
在黑板上画数轴，标出 $0,2,4,6,8,10$。画一条正弦式的连续波浪线，令 $x=0, 1, 2$ 为零点。\par
让学生用红笔在数轴上逐个圈出零点：
\begin{itemize}[leftmargin=1em, itemsep=0.1em]
  \item $[0,2]$: $0, 1, 2$ （3 个）
  \item $(2,4]$: $3, 4$ （2 个）
  \item $(4,6]$: $5, 6$ （2 个）
  \item $(6,8]$: $7, 8$ （2 个）
  \item $(8,10]$: $9, 10$ （2 个）
\end{itemize}
总计 $3 + 2 + 2 + 2 + 2 = 13$ 个零点。
}
\end{paracol}

\end{document}
